
\section{Methods}
We will introduce a range of very practical and simple to implement ideas towards our goal of leveraging richer preference feedback during fine-tuning.

\subsection{Idea 1: Jointly maximise the rationale }


Consider expanding the DPO objective as follows:

\begin{equation}
    \mathcal{L}_{\theta}(\mathcal{X}_{R}) = - \mathbb{E}_{\mathcal{X}_{R}}
    \left[ \log \sigma
    \left(
    \hat{r}(x, y_{w})
    - \hat{r}(x, y_{l})
    \right)
    + \log p_{\theta}(r | c :y_{w}, y_{l})
    \right],
    \label{eq:cdpo_loss}
\end{equation}

Where $\mathcal{X}_{R}$ is our preference dataset expanded with the reasoning for the given preference judgments: $\mathcal{X}_{R} = \{x, r, y_{w}, y_{l}\}^{N}$ and $c$ are the instructions provided to the oracle for providing the reasoning. $c :y_{w}, y_{l}$ indicates that these instructions are combined with the preference labels as context.

\subsubsection{Experiments}


\paragraph{}
For code, given the simplicity of the change, we should explore
\url{https://huggingface.co/docs/trl/main/en/dpo_trainer}

\paragraph{}
We can also reference the Active Preference Learning code:
\url{https://github.com/willmuldrew/active-preference-learning}



\subsection{Idea 2: Bayes rule over rationale}

Consider the following formulation of our language model where we explicitly model the rationale:

\begin{equation}
    p_{\theta}(y_{w} \succ y_{l}, r | x) = p_{\theta}(y_{w} \succ y_{l} | x) p_{\theta}(r | x, y_{w} \succ y_{l} )
\end{equation}

What about treating it as a latent variable and creating a variational objective?


\subsection{Idea 3: Corrections from feedback}

We have a model $p_\theta(r|q)$. Here $q = q_1,\ldots, q_{Q}$ is the query context and we want to generate an output response $r_1,\ldots, r_R$.

In a standard training setup we would maximise the likelihood of context, generation pairs $\sum_n \log p_\theta(r^n|q^n)$.\\

Let's imagine we have though something like a setup whereby a user can provide a response $u$ to the generation, along with a positive/negative reward $y\in\{1,-1\}$.\\

For example,
\begin{align}
    q & =\text{what is the capital of France?}  \\
    r & =\text{London is the capital of France} \\
    y & =-1                                     \\
    u & =\text{The capital of France is Paris}
\end{align}
In this case we want the model to decrease the probability of $\log p_\theta(r|q)$ but to increase the probability of the correct response $\log p_\theta(u|q)$.  In this case it seems straightforward.  One objective would be to increase  $\log \sigma\left(\log p_\theta(u|q) - \log p_\theta(r|q)\right)$.\\

The above example is straightforward. A more complex situation is when for example:
\begin{align}
    q & =\text{Write a four line poem}                               \\
    r & =\text{Mary had a little lamb; it's feet were white as snow} \\
    y & =-1                                                          \\
    u & =\text{That's only a two line poem}
\end{align}
In this case we don't have a correct answer for the model. One approach would be to use an LLM to write a correct question-answer pair, based on the user correction. See the image:
\begin{center}

\end{center}
In this case GPT learns the correct query-answer pair in both cases. One then wants to decrease the probability of the original response and increase the probability of the correct query-response $q',r'$. One could then add a term like
\begin{equation}
    \log \sigma\left(\log p_\theta(r'|q') - \log p_\theta(r|q) \right)
\end{equation}

One may of course wish to normalise these probabilities by their respective sequence lengths.

I don't know how good GPT is at generating correct query-answer pairs like this, but if it's quite good then this may be a simple way to make use of richer feedback to fine-tune a model.